\chapter{Package and Pinout}
\label{ch:package}

The video device is a single integrated circuit with its video memory in separate static RAM chips. The picture
processing unit, the palette, the sprite table and the colour converters are on the one piece of silicon. The 64 KB of
video memory is outside it, on a memory bus of its own, so besides the bus of the CPU, two interrupt outputs, the clock
and reset inputs and the video output, the package carries the signals of that memory bus. This chapter gives the
package, the meaning of every pin and the timing the video memory must meet. Chapter~\ref{ch:signal-timing} gives the
timing of the signals.

\section{Chip Layout}

The device is supplied in a 100-pin plastic quad flat package (QFP-100), 25 pins on each side. Pin 1 is at the top left
corner, marked by a dot on the package; the pins are numbered counter-clockwise when the package is seen from above.

\begin{figure}[H]
	\centering
	\begin{tikzpicture}[
			scale=0.74, transform shape,
			pinlabel/.style={font=\tiny},
			pinnum/.style={font=\tiny, text=sircgray},
			body/.style={draw, thick, fill=gray!8},
			pin/.style={draw, thick}
		]
		\def\pitch{0.4}
		\def\side{10.6}
		\def\first{0.5}
		% Package body
		\draw[body] (0,0) rectangle (\side,\side);
		% Pin 1 index dot
		\fill (0.55,\side-0.55) circle (0.08);
		% Left side, top to bottom: pins 1 to 25
		\foreach \n/\name [count=\i from 0] in {
				1/VCC1, 2/CLKI, 3/\activelow{RST}, 4/\activelow{CS}, 5/\activelow{BAS}, 6/BRW, 7/\activelow{BACK}, 8/A0,
				9/A1, 10/A2, 11/A3, 12/A4, 13/A5, 14/A6, 15/A7, 16/A8,
				17/A9, 18/A10, 19/A11, 20/A12, 21/A13, 22/A14, 23/A15, 24/GND1,
				25/VCC2
			} {
				\pgfmathsetmacro{\y}{\side - \first - \i*\pitch}
				\draw[pin] (0,\y) -- (-0.4,\y);
				\node[pinnum, anchor=west] at (0.1,\y) {\n};
				\node[pinlabel, anchor=east] at (-0.46,\y) {\name};
			}
		% Bottom side, left to right: pins 26 to 50
		\foreach \n/\name [count=\i from 0] in {
				26/D0, 27/D1, 28/D2, 29/D3, 30/D4, 31/D5, 32/D6, 33/D7,
				34/D8, 35/D9, 36/D10, 37/D11, 38/D12, 39/D13, 40/D14, 41/D15,
				42/GND2, 43/\activelow{NMI}, 44/\activelow{IRQ4}, 45/\activelow{CSYNC}, 46/AVCC, 47/IREF, 48/R, 49/G,
				50/B
			} {
				\pgfmathsetmacro{\x}{\first + \i*\pitch}
				\draw[pin] (\x,0) -- (\x,-0.4);
				\node[pinnum, rotate=90, anchor=west] at (\x,0.1) {\n};
				\node[pinlabel, rotate=90, anchor=east] at (\x,-0.46) {\name};
			}
		% Right side, bottom to top: pins 51 to 75
		\foreach \n/\name [count=\i from 0] in {
				51/AGND, 52/VCC3, 53/VA0, 54/VA1, 55/VA2, 56/VA3, 57/VA4, 58/VA5,
				59/VA6, 60/VA7, 61/VA8, 62/VA9, 63/VA10, 64/VA11, 65/VA12, 66/VA13,
				67/VA14, 68/GND3, 69/VCC4, 70/VD0, 71/VD1, 72/VD2, 73/VD3, 74/VD4,
				75/VD5
			} {
				\pgfmathsetmacro{\y}{\first + \i*\pitch}
				\draw[pin] (\side,\y) -- (\side+0.4,\y);
				\node[pinnum, anchor=east] at (\side-0.1,\y) {\n};
				\node[pinlabel, anchor=west] at (\side+0.46,\y) {\name};
			}
		% Top side, right to left: pins 76 to 100
		\foreach \n/\name [count=\i from 0] in {
				76/VD6, 77/VD7, 78/VD8, 79/VD9, 80/VD10, 81/VD11, 82/VD12, 83/VD13,
				84/VD14, 85/VD15, 86/GND4, 87/\activelow{VCE}, 88/\activelow{VOE}, 89/\activelow{VWE}, 90/NC, 91/NC,
				92/NC, 93/NC, 94/NC, 95/NC, 96/NC, 97/NC, 98/NC, 99/NC,
				100/NC
			} {
				\pgfmathsetmacro{\x}{\side - \first - \i*\pitch}
				\draw[pin] (\x,\side) -- (\x,\side+0.4);
				\node[pinnum, rotate=90, anchor=east] at (\x,\side-0.1) {\n};
				\node[pinlabel, rotate=90, anchor=west] at (\x,\side+0.46) {\name};
			}
		% Package label
		\node[font=\bfseries] at (\side/2,\side/2+0.25) {SIRCADA};
		\node[font=\small] at (\side/2,\side/2-0.25) {100-pin QFP};
	\end{tikzpicture}
	\caption{SIRCADA 100-pin QFP Pinout, Top View}
	\label{fig:pinout}
\end{figure}

\begin{table}[H]
	\centering
	\begin{tabular}{rll}
		\toprule
		\textbf{Pin} & \textbf{Name}     & \textbf{Description}              \\
		\midrule
		1            & VCC1              & Power supply                      \\
		2            & CLKI              & Clock input                       \\
		3            & \activelow{RST}   & Reset input                       \\
		4            & \activelow{CS}    & Chip select input                 \\
		5            & \activelow{BAS}   & Bus access strobe input           \\
		6            & BRW               & Bus read/write input              \\
		7            & \activelow{BACK}  & Bus acknowledge output            \\
		8--23        & A0--A15           & Address bus bits 0--15            \\
		24           & GND1              & Ground                            \\
		25           & VCC2              & Power supply                      \\
		26--41       & D0--D15           & Data bus bits 0--15               \\
		42           & GND2              & Ground                            \\
		43           & \activelow{NMI}   & Vertical sync interrupt output    \\
		44           & \activelow{IRQ4}  & Line-compare interrupt output     \\
		45           & \activelow{CSYNC} & Composite sync output             \\
		46           & AVCC              & Analog power supply               \\
		47           & IREF              & Colour output reference           \\
		48--50       & R, G, B           & Analog colour outputs             \\
		51           & AGND              & Analog ground                     \\
		52           & VCC3              & Power supply                      \\
		53--67       & VA0--VA14         & Video memory address bits 0--14   \\
		68           & GND3              & Ground                            \\
		69           & VCC4              & Power supply                      \\
		70--85       & VD0--VD15         & Video memory data bits 0--15      \\
		86           & GND4              & Ground                            \\
		87           & \activelow{VCE}   & Video memory chip enable output   \\
		88           & \activelow{VOE}   & Video memory output enable output \\
		89           & \activelow{VWE}   & Video memory write enable output  \\
		90--100      & NC                & No connection                     \\
		\bottomrule
	\end{tabular}
	\caption{SIRCADA 100-pin QFP Pinout}
	\label{tab:pinout}
\end{table}

\section{Signal Descriptions}

Table~\ref{tab:signals} lists the direction and sense of each signal. Active-low signals carry a bar over the name;
they are asserted when low. Address and data signals are active high. An output marked \emph{open drain} can only pull
its line low; the system supplies a pull-up, and several devices may share the line.

\begin{table}[H]
	\centering
	\small
	\begin{tabularx}{\textwidth}{llllX}
		\toprule
		\textbf{Signal}   & \textbf{Pins} & \textbf{Direction} & \textbf{Output type} & \textbf{Function}                             \\
		\midrule
		CLKI              & 2             & Input              &                      & Rising edge advances the device one PPU cycle \\
		\activelow{RST}   & 3             & Input              &                      & Resets the device                             \\
		\activelow{CS}    & 4             & Input              &                      & Selects the device for a bus cycle            \\
		\activelow{BAS}   & 5             & Input              &                      & Marks a bus cycle                             \\
		BRW               & 6             & Input              &                      & High for a read, low for a write              \\
		\activelow{BACK}  & 7             & Output             & Open drain           & Completes the bus cycle                       \\
		A0--A15           & 8--23         & Input              &                      & Word offset in the device segment             \\
		D0--D15           & 26--41        & Bidirectional      & Three-state          & Data                                          \\
		\activelow{NMI}   & 43            & Output             & Open drain           & Vertical sync interrupt                       \\
		VA0--VA14         & 53--67        & Output             & Driven               & Video memory word address                     \\
		VD0--VD15         & 70--85        & Bidirectional      & Three-state          & Video memory data                             \\
		\activelow{VCE}   & 87            & Output             & Driven               & Enables the video memory                      \\
		\activelow{VOE}   & 88            & Output             & Driven               & Video memory drives its data                  \\
		\activelow{VWE}   & 89            & Output             & Driven               & Video memory takes a word                     \\
		\activelow{IRQ4}  & 44            & Output             & Open drain           & Line-compare interrupt                        \\
		\activelow{CSYNC} & 45            & Output             & Driven               & Composite sync                                \\
		R, G, B           & 48--50        & Output             & Analog               & Colour                                        \\
		IREF              & 47            & Reference          &                      & Sets the full-scale colour level              \\
		\bottomrule
	\end{tabularx}
	\caption{Signal Summary}
	\label{tab:signals}
\end{table}

\subsection{Clock Input (CLKI)}

The device has one clock input, the same \texttt{CLKI} the system supplies to the CPU. One \texttt{CLKI} cycle is one
PPU cycle: the master clock divided by two (Section~\ref{sec:cpu-clock}). Every internal interval, the pixel, the group
of four pixels, the line and the frame, is a whole number of \texttt{CLKI} cycles, so the device needs no other clock.
The pixel rate is \texttt{CLKI} divided by two. The rising edge is the only edge the device uses. The range of clock
frequencies is not yet specified.

\subsection{Reset Input (\activelow{RST})}

While \activelow{RST} is asserted the device is held in its reset state (Section~\ref{sec:signal-reset}). The input is
sampled on each rising edge of \texttt{CLKI}, and must stay asserted for at least six cycles. The \activelow{RSTO}
output of the CPU holds for at least that long and can be connected directly.

\subsection{Bus Signals}

The device is a bus peripheral. It watches the bus of the CPU (\hwmanual{}) and answers the cycles that are addressed
to it. It has no byte strobes and no access-type inputs: it treats an instruction fetch, a data access and a DMA
transfer alike.

\begin{description}
	\item[\activelow{CS}] The system decodes the segment of the device, the upper address bits that the device does not see, and
	      asserts \activelow{CS} while the segment is addressed. A cycle in which \activelow{CS} or \activelow{BAS} is
	      deasserted is ignored: the device drives neither the data bus nor \activelow{BACK}.
	\item[\activelow{BAS}, BRW] The CPU asserts \activelow{BAS} for a whole cycle that makes a bus access. BRW is high for a
	      read and low for a write. Both are valid for the whole of the cycle.
	\item[A0--A15] The word offset of the access within the segment (Chapter~\ref{ch:memory-map}). The device accesses one 16-bit
	      word per cycle.
	\item[D0--D15] On a write the CPU drives the bus and the device takes the word at the rising edge that ends the cycle. On a
	      read the device drives the word that the offset selects, from the start of the cycle to its end, and releases the bus
	      when the cycle ends. A read of an offset that has no register or memory is driven as zero.
	\item[\activelow{BACK}] The device asserts \activelow{BACK} for the whole of every cycle in which \activelow{CS} and \activelow{BAS} are asserted, whatever the offset and whatever the state of the picture, with one exception: for a read of video memory it holds \activelow{BACK} off for the first cycle, which makes the CPU wait one cycle (Section~\ref{sec:signal-vram}), or two when a write to video memory immediately precedes it. It never raises a bus error. \activelow{BACK} is open drain, so it may be joined with the acknowledge from other devices on the bus.
\end{description}

\subsection{Video Memory Bus (VA0--VA14, VD0--VD15, \activelow{VCE}, \activelow{VOE}, \activelow{VWE})}
\label{sec:package-vram-bus}

The video memory is 32,768 words of 16 bits, for instance two static RAMs of 32,768 by 8 bits, one for the low byte and
one for the high byte of every word. The device is the only master of this bus. The CPU reaches video memory through
the device, in the cycles the renderer leaves free (Section~\ref{sec:memory-vram-access}).

\begin{description}
	\item[VA0--VA14] The word address of the access. Address bit 15 of the offset in the segment (\texttt{0x8000}) selects video
	      memory and is not part of this address.
	\item[VD0--VD15] On a read the memory drives the word and the device takes it at the rising edge of \texttt{CLKI} that ends
	      the cycle. On a write the device drives the word for the whole of the cycle in which \activelow{VWE} is asserted.
	\item[\activelow{VCE}, \activelow{VOE}, \activelow{VWE}] \activelow{VCE} is asserted in every cycle in which the device
	      accesses the memory. \activelow{VOE} is asserted as well for a read and \activelow{VWE} for a write.
\end{description}

The memory must meet the requirements of Table~\ref{tab:vram-timing}, which are those of one access in every cycle. A
cycle is one \texttt{CLKI} cycle, 93~ns at the 10.74~MHz of the SIRCUS system. The values are maxima, measured at the
pins of the device, so the delays of the board and of the pins of the device are already counted in them.

\begin{table}[H]
	\centering
	\begin{tabular}{lr}
		\toprule
		\textbf{Requirement}                          & \textbf{Maximum} \\
		\midrule
		Address valid to data valid (read)            & 70~ns            \\
		\activelow{VCE} asserted to data valid (read) & 70~ns            \\
		\activelow{VOE} asserted to data valid (read) & 70~ns            \\
		Address valid to end of write                 & 70~ns            \\
		\activelow{VCE} asserted to end of write      & 70~ns            \\
		Write pulse width                             & 70~ns            \\
		Data valid to end of write                    & 70~ns            \\
		Write cycle time                              & 93~ns            \\
		\bottomrule
	\end{tabular}
	\caption{Video Memory Timing Requirements}
	\label{tab:vram-timing}
\end{table}

The memory must need no address or data hold time after the end of a write. Which memory parts meet the requirements,
and how they are placed on the board, is for the system to decide.

\subsection{Interrupt Outputs (\activelow{NMI}, \activelow{IRQ4})}

The vertical sync interrupt is a one-cycle pulse on \activelow{NMI}, and the line-compare interrupt is a level held on
\activelow{IRQ4} (Section~\ref{sec:interrupts}). Both are open drain, so the line is shared with other devices when the
system connects them; a program that shares \activelow{IRQ4} reads the \texttt{LIP} bit of the Status register to see
whose interrupt it is. Neither output is driven while \activelow{RST} is asserted.

\subsection{Video Output}

\begin{description}
	\item[\activelow{CSYNC}] Composite sync. It is asserted for the horizontal sync pulse of every line and, in the lines of
	      vertical sync, for the vertical sync pulses (Section~\ref{sec:video-output}).
	\item[R, G, B] The picture, as three analog levels. Each is converted from 5 bits of the colour of the pixel after the
	      brightness setting has been applied, so it has 32 steps from black to full scale. In blanking, in forced blank and
	      during reset the three outputs are at black.
	\item[IREF] A resistor from IREF to AGND sets the level of full scale for all three outputs. The value of the resistor and
	      the voltages of the outputs are not yet specified.
	\item[AVCC, AGND] The supply of the converters, kept apart from VCC1--VCC4 and GND1--GND4 so that the noise of the logic does
	      not reach the picture. AVCC is 5~V; AGND is connected to the ground of the logic at one point on the board.
\end{description}

\subsection{Power Pins (VCC1--VCC4, GND1--GND4)}

The eight pins supply the logic and the pads of the video memory bus, and are all connected. VCC is 5~V $\pm$ 5\%, as
for the CPU. The electrical characteristics other than the supply are not yet specified.

\section{Connecting the Device}

Figure~\ref{fig:connections} shows the device in a system. The CPU and the video device share the address bus (the
device uses the lower 16 bits), the data bus, \activelow{BAS}, BRW, \activelow{BACK} and \texttt{CLKI}. The address
decoder asserts \activelow{CS} for the segment of the device. The reset output of the CPU resets the device. The video
memory is joined to the memory bus of the device and to nothing else. The interrupt outputs of the device go to the
interrupt inputs of the CPU, with pull-up resistors, and the video outputs go to the display connector.

\begin{figure}[H]
	\centering
	\begin{tikzpicture}[
			scale=0.9, transform shape,
			block/.style={draw, rounded corners=2pt, fill=sirclightgray, align=center, font=\small},
			sig/.style={-{Stealth[length=2mm]}, thick},
			bus/.style={{Stealth[length=2mm]}-{Stealth[length=2mm]}, ultra thick},
			note/.style={font=\scriptsize, text=sircgray, align=center}
		]
		\node[block, minimum width=2.4cm, minimum height=6.0cm] (cpu) at (0,0) {SIRC-1\\CPU};
		\node[block, minimum width=2.6cm, minimum height=6.0cm] (ppu) at (8.0,0) {SIRCADA};
		\node[block, minimum width=2.2cm, minimum height=1.0cm] (dec) at (4.0,4.1) {Address\\decode};
		\node[block, minimum width=2.2cm, minimum height=1.0cm] (clk) at (4.0,-4.1) {Clock\\generator};
		\node[block, minimum width=2.6cm, minimum height=1.0cm] (out) at (12.4,0) {Display\\connector};
		\node[block, minimum width=2.6cm, minimum height=1.0cm] (vram) at (12.4,-4.1) {Video memory\\32K by 16};
		% Lanes between the CPU and the device
		\draw[bus] (1.2,2.0) -- node[above, note]{A0--A15} (6.7,2.0);
		\draw[bus] (1.2,1.2) -- node[above, note]{D0--D15} (6.7,1.2);
		\draw[sig] (1.2,0.4) -- node[above, note]{\activelow{BAS}, BRW} (6.7,0.4);
		\draw[sig] (6.7,-0.4) -- node[above, note]{\activelow{BACK}} (1.2,-0.4);
		\draw[sig] (1.2,-1.2) -- node[above, note]{\activelow{RSTO} to \activelow{RST}} (6.7,-1.2);
		\draw[sig] (6.7,-2.0) -- node[above, note]{\activelow{NMI}, \activelow{IRQ4} (pull-ups)} (1.2,-2.0);
		% Address decode
		\draw[sig] (cpu.north) |- node[above, pos=0.75, note]{A16--A23} (dec.west);
		\draw[sig] (dec.east) -| node[above, pos=0.25, note]{\activelow{CS}} (ppu.north);
		% Clock
		\draw[sig] (clk.west) -| node[below, pos=0.25, note]{CLKI} (cpu.south);
		\draw[sig] (clk.east) -| (ppu.south);
		% Video
		\draw[sig] (ppu.east) -- node[above, note]{R, G, B,\\\activelow{CSYNC}} (out.west);
		% Video memory bus
		\draw[bus] ([xshift=0.8cm]ppu.south) |- node[below, pos=0.55, note]{VA, VD, \activelow{VCE}, \activelow{VOE}, \activelow{VWE}} (vram.west);
	\end{tikzpicture}
	\caption{The Video Device in a System}
	\label{fig:connections}
\end{figure}
